\documentclass[11pt]{article}

% ---------- Packages ----------
\usepackage[margin=1in]{geometry}
\usepackage{amsmath,amssymb}
\usepackage{booktabs}
\usepackage{array}
\usepackage{enumitem}
\usepackage{xcolor}
\usepackage[most]{tcolorbox}
\usepackage{titlesec}
\usepackage{parskip}
\usepackage{hyperref}
\usepackage{fancyhdr}

% ---------- Hyperref ----------
\hypersetup{
  colorlinks=true,
  linkcolor=blue!50!black,
  urlcolor=blue!50!black,
  pdftitle={AP Statistics: AI Tutoring, Compliance, and Evidence}
}

% ---------- Header / footer ----------
\pagestyle{fancy}
\fancyhf{}
\fancyhead[L]{\small AP Statistics --- AI Tutoring, Compliance, and Evidence}
\fancyhead[R]{\small \thepage}
\renewcommand{\headrulewidth}{0.4pt}

% ---------- Section formatting ----------
\titleformat{\part}[display]
  {\normalfont\Huge\bfseries\centering}
  {\partname\ \thepart}{0.5em}{}
\titleformat{\section}
  {\normalfont\Large\bfseries\color{blue!50!black}}{\thesection}{1em}{}
\titleformat{\subsection}
  {\normalfont\large\bfseries\color{blue!40!black}}{\thesubsection}{1em}{}

% ---------- Boxes ----------
\newtcolorbox{scenario}{
  colback=gray!5, colframe=gray!50,
  title=Scenario, fonttitle=\bfseries,
  boxrule=0.5pt, arc=2pt
}
\newtcolorbox{notebox}{
  colback=blue!3, colframe=blue!40!black,
  title=Note, fonttitle=\bfseries,
  boxrule=0.5pt, arc=2pt
}
\newtcolorbox{expectedbox}{
  colback=green!3, colframe=green!40!black,
  title=Expected Response, fonttitle=\bfseries,
  boxrule=0.5pt, arc=2pt
}
\newtcolorbox{keybox}{
  colback=yellow!8, colframe=orange!60!black,
  title=Worked Answer Key, fonttitle=\bfseries,
  boxrule=0.5pt, arc=2pt
}

% ---------- Custom commands ----------
\newcommand{\code}[1]{\texttt{#1}}
\newcommand{\phat}{\hat{p}}
\newcommand{\xbar}{\bar{x}}

\title{\bfseries AP Statistics:\\ AI Tutoring, Compliance, and Evidence \\[0.4em]
\large A Combined Capstone and Investigative Task}
\author{}
\date{}

\begin{document}
\maketitle
\thispagestyle{empty}

\begin{center}
\begin{minipage}{0.85\textwidth}
\itshape
This document combines two complementary AP Statistics investigations.
\textbf{Part I} is a 60-question synthetic capstone covering Units 1--9,
designed to assess every major inferential and conceptual skill in the
course. \textbf{Part II} is a focused investigative task that applies
those same tools to a real-world causal question about AI-tutor
pedagogy, drawing on a published 2025 randomized controlled trial.
Together they distinguish \emph{intent-to-treat effects},
\emph{on-treatment associations}, \emph{confounding}, and
\emph{selection bias}.
\end{minipage}
\end{center}

\vspace{1em}
\tableofcontents
\newpage

% ============================================================
\part{Capstone: AI Tutoring, Compliance, and Evidence}
% ============================================================

\begin{notebox}
This is a synthetic AP Statistics investigation. The scenario is
realistic, but the data are invented so the statistical reasoning can
be evaluated cleanly.
\end{notebox}

\section{Scenario}

A district is debating whether AI tutoring should be allowed during the
learning phase of AP Statistics. Some teachers argue that AI mainly
enables cheating. Others argue that AI works like a private tutor:
useful during practice, but not a substitute for independent reasoning
on assessments.

The district runs a six-week study with 240 AP Statistics students
from 8 schools. Within each school, students are randomly assigned to
one of two instructional conditions:

\begin{itemize}[leftmargin=*]
  \item \textbf{AI-access condition:} Students may use an AI tutor
  during practice, homework, and review. They must still show
  handwritten work and complete in-class checks without AI.
  \item \textbf{Standard-resource condition:} Students use AP
  Classroom videos, teacher notes, textbook examples, and ordinary
  peer discussion. They may not use AI during the study period.
\end{itemize}

Researchers measure each student's pretest score, posttest score,
daily AI-tutor minutes, phone-policy environment, engagement rating,
and whether the student passes a later AP-style practice exam. Use a
significance level of $\alpha = 0.05$ unless otherwise stated.

\section{Data}

\subsection*{Table A: One-Variable Data}
For 20 randomly selected students in the AI-access condition, the
recorded average daily AI-tutor minutes were:
\[
0,\ 0,\ 5,\ 8,\ 10,\ 12,\ 15,\ 18,\ 22,\ 25,\ 30,\ 32,\ 35,\ 40,\ 45,\ 50,\ 55,\ 60,\ 75,\ 95
\]

\subsection*{Table B: Randomized Study Summary}
\begin{center}
\small
\begin{tabular}{lrrrrrrrr}
\toprule
Condition & $n$ & Mean pre & SD pre & Mean post & SD post & Mean gain & SD gain & Passed \\
\midrule
AI-access          & 120 & 48.2 & 13.8 & 66.8 & 15.6 & 18.6 & 12.4 & 69 \\
Standard-resource  & 120 & 48.9 & 14.1 & 58.7 & 16.2 &  9.8 & 11.7 & 47 \\
\bottomrule
\end{tabular}
\end{center}

\subsection*{Table C: Actual AI Use Within the AI-Access Condition}
\begin{center}
\small
\begin{tabular}{lrrrr}
\toprule
Actual behavior & $n$ & Mean gain & SD gain & Passed \\
\midrule
Regular use ($\geq 20$ min/day)        & 76 & 23.4 & 11.6 & 55 \\
Little or no use ($< 20$ min/day)      & 44 & 10.3 & 10.9 & 14 \\
\bottomrule
\end{tabular}
\end{center}

\subsection*{Table D: School Culture and AI Use}
\begin{center}
\small
\begin{tabular}{lrrrr}
\toprule
School culture & Regular & Occasional & None & Total \\
\midrule
Restrictive phone/AI culture & 21 & 39 & 60 & 120 \\
Open AI-tutor culture        & 66 & 34 & 20 & 120 \\
\midrule
\textbf{Total}               & \textbf{87} & \textbf{73} & \textbf{80} & \textbf{240} \\
\bottomrule
\end{tabular}
\end{center}

\subsection*{Table E: Regression Summary}
For 80 students with complete usage logs, researchers fit a
least-squares regression model predicting gain score from average
daily AI-tutor minutes:
\[
\widehat{\text{gain}} = 6.2 + 0.085(\text{minutes})
\]
Additional information:
\begin{itemize}[leftmargin=*,nosep]
  \item $n = 80$
  \item $r = 0.62$
  \item $s = 9.8$
  \item $SE_b = 0.013$ (standard error of the slope)
  \item Residual plot shows random scatter with roughly constant vertical spread.
  \item Histogram of residuals is roughly unimodal and approximately symmetric.
\end{itemize}

\subsection*{Table F: Probability Model}
Based on a larger district survey, suppose an AP Statistics student
independently has probability $p = 0.32$ of becoming a regular
AI-tutor user when AI access is allowed but not required.

\section{Questions}

\subsection*{Unit 1: Exploring One-Variable Data}
\begin{enumerate}[leftmargin=*]
\item Classify each variable as categorical or quantitative:
instructional condition, school culture, average daily AI-tutor
minutes, gain score, passed AP-style practice exam.
\item Use Table A to make an appropriate graph of AI-tutor minutes.
Describe shape, center, spread, and unusual features in context.
\item Find the median, $Q_1$, $Q_3$, IQR, and range of Table A.
\item Use the $1.5 \times \text{IQR}$ rule to determine whether 95
minutes is an outlier.
\item Explain why the mean AI-tutor minutes would be larger than the
median for Table A.
\end{enumerate}

\subsection*{Unit 2: Exploring Two-Variable Data}
\begin{enumerate}[resume,leftmargin=*]
\item Use Table D to compute the conditional relative frequency of
regular AI use within each school culture.
\item Does Table D suggest an association between school culture and
actual AI use? Explain using conditional relative frequencies.
\item Interpret the slope $0.085$ in Table E in context.
\item Interpret $r = 0.62$ in context.
\item Find and interpret $r^2$.
\item Predict the gain score for a student who uses the AI tutor for
40 minutes per day.
\item Explain why it would be risky to use the regression line to
predict the gain score for a student who uses AI for 180 minutes
per day.
\end{enumerate}

\subsection*{Unit 3: Collecting Data}
\begin{enumerate}[resume,leftmargin=*]
\item Identify the experimental units, treatments, and response
variables in the six-week study.
\item Explain why random assignment within each school improves the
study.
\item Explain why blocking by school would be sensible in this study.
\item Can the randomized comparison in Table B support a
cause-and-effect conclusion about AI access? Explain.
\item Can the comparison in Table C support a cause-and-effect
conclusion about actual AI use? Explain carefully.
\item Identify one possible source of bias if students self-report
their AI minutes.
\item Explain how noncompliance can contaminate the interpretation of
a treatment effect.
\end{enumerate}

\subsection*{Unit 4: Probability, Random Variables, and Probability Distributions}
\begin{enumerate}[resume,leftmargin=*]
\item Using Table D, find $P(\text{regular AI use})$.
\item Using Table D, find $P(\text{open AI-tutor culture})$.
\item Using Table D, find $P(\text{regular AI use} \mid \text{open AI-tutor culture})$.
\item Are regular AI use and open AI-tutor culture independent in
Table D? Justify numerically.
\item Suppose 5 students are randomly selected from a population
where $p = 0.32$ for regular AI-tutor use. Find the probability
exactly 2 are regular AI-tutor users.
\item Let $X$ be the number of regular AI-tutor users among 5
randomly selected students with $p = 0.32$. Find $E(X)$ and
$\text{SD}(X)$.
\item Find the probability that the first regular AI-tutor user
appears on the fourth student selected.
\end{enumerate}

\subsection*{Unit 5: Sampling Distributions}
\begin{enumerate}[resume,leftmargin=*]
\item For samples of size 120 from a population with true pass
proportion $p = 0.50$, describe the approximate sampling
distribution of $\hat{p}$. Check conditions.
\item If the true mean gain under the standard-resource condition is
10 points with population SD 12, describe the approximate sampling
distribution of the sample mean gain for $n = 120$.
\item Under the null hypothesis that the two pass proportions are
equal, explain why the observed difference in pass rates from
Table B might or might not be surprising.
\item Explain why larger sample sizes make it easier to detect
smaller real differences.
\end{enumerate}

\subsection*{Unit 6: Inference for Categorical Data: Proportions}
\begin{enumerate}[resume,leftmargin=*]
\item Construct a 95\% confidence interval for the proportion of
AI-access students who passed the AP-style practice exam.
\item Construct a 95\% confidence interval for the difference in pass
proportions, $p_{AI} - p_{\text{standard}}$, using Table B.
\item Interpret the interval from Question 32 in context.
\item Test whether the AI-access condition produced a higher pass
proportion than the standard-resource condition.
\item Explain the meaning of the p-value from Question 34 in context.
\item Describe a Type I error and a Type II error for the test in
Question 34.
\end{enumerate}

\subsection*{Unit 7: Inference for Quantitative Data: Means}
\begin{enumerate}[resume,leftmargin=*]
\item Construct a 95\% confidence interval for the mean gain in the
AI-access condition.
\item Construct a 95\% confidence interval for the difference in mean
gains, $\mu_{AI} - \mu_{\text{standard}}$, using Table B.
\item Interpret the interval from Question 38 in context.
\item Test whether the AI-access condition produced a greater mean
gain than the standard-resource condition.
\item Explain why the randomized design matters when interpreting the
difference in mean gains.
\item Suppose a critic says, ``Students in the AI-access group were
already stronger.'' Use Table B to respond statistically.
\end{enumerate}

\subsection*{Unit 8: Inference for Categorical Data: Chi-Square}
\begin{enumerate}[resume,leftmargin=*]
\item Use Table D to state appropriate hypotheses for a chi-square test.
\item Calculate the expected count for restrictive culture and
regular AI use.
\item Verify that the conditions for a chi-square test are met.
\item The chi-square test statistic for Table D is approximately
$43.62$. Find the degrees of freedom and approximate p-value
category.
\item State a conclusion in context.
\item Identify which cells contribute most to the chi-square
statistic and explain what that means in context.
\end{enumerate}

\subsection*{Unit 9: Inference for Quantitative Data: Slopes}
\begin{enumerate}[resume,leftmargin=*]
\item State hypotheses for testing whether average daily AI-tutor
minutes are linearly associated with gain score.
\item Use Table E to calculate the test statistic for the slope.
\item Find the degrees of freedom.
\item State a conclusion in context.
\item Construct a 95\% confidence interval for the slope.
\item Interpret the interval in context.
\item Explain why regression evidence alone is weaker causal evidence
than the randomized comparison in Table B.
\end{enumerate}

\subsection*{Final Synthesis}
\begin{enumerate}[resume,leftmargin=*]
\item A school administrator says, ``The AP-style scores prove
AI-supported instruction works.'' Explain why that statement is
too strong.
\item A teacher says, ``The AP-style scores prove AI-supported
instruction fails whenever students do not use it correctly.''
Explain why that statement is also too strong.
\item Distinguish between the \textbf{intent-to-treat effect} of AI
access and the \textbf{on-treatment association} for students who
actually used AI regularly.
\item Write a statistically careful conclusion about what the full
dataset does and does not show.
\item Propose a better follow-up study that could distinguish
AI-tutor effects from student motivation, compliance, teacher
culture, and phone-policy effects.
\end{enumerate}

\newpage
\section{Teacher Notes / Answer Key}

\begin{keybox}
Worked answers below. Approximations are AP-appropriate; minor
rounding differences are acceptable.
\end{keybox}

\begin{enumerate}[leftmargin=*]
\item Condition, school culture, and passed exam are categorical.
AI minutes and gain score are quantitative.

\item A dotplot, histogram, or boxplot is appropriate. The
distribution is right-skewed with high values at 75 and 95
minutes. A typical value is around 25 to 30 minutes. The data
range from 0 to 95 minutes.

\item Median $27.5$; $Q_1 = 11$; $Q_3 = 47.5$; IQR $= 36.5$; range $= 95$.

\item Lower fence $11 - 1.5(36.5) = -43.75$; upper fence
$47.5 + 1.5(36.5) = 102.25$. The value 95 is \emph{not} an outlier
by the $1.5 \times \text{IQR}$ rule.

\item The large high-end values create a right tail, and the mean is
pulled toward that tail.

\item Restrictive regular-use rate: $21/120 = 0.175$. Open-culture
regular-use rate: $66/120 = 0.55$.

\item Yes. Regular use is much higher in open AI-tutor culture, while
no use is much higher in restrictive culture: $60/120 = 0.50$
versus $20/120 = 0.167$.

\item For each additional AI-tutor minute per day, predicted gain
increases by $0.085$ points, on average.

\item There is a moderately strong positive linear association
between AI-tutor minutes and gain score.

\item $r^2 = 0.62^2 = 0.3844$. About 38.4\% of variation in gain
score is explained by the linear relationship with AI-tutor
minutes.

\item $6.2 + 0.085(40) = 9.6$ points.

\item 180 minutes is likely outside the observed range, so this is
extrapolation. The linear pattern may not continue that far.

\item Experimental units are AP Statistics students. Treatments are
AI access and standard resources. Response variables include
posttest score, gain score, pass status, and engagement.

\item Random assignment reduces systematic differences between
groups, making the groups more comparable except for the assigned
condition.

\item Blocking by school controls for school-level differences such
as teacher culture, phone policy, scheduling, prior preparation,
and local norms.

\item Yes, the randomized comparison can support a cause-and-effect
conclusion about \emph{being assigned} AI access, assuming the
study was carried out properly.

\item No, not by itself. Actual use was not randomly assigned; it may
be confounded with motivation, prior skill, family support,
teacher encouragement, or school culture.

\item Students may overreport use to look diligent, underreport use
if they fear punishment, or misremember their actual minutes.

\item Noncompliance means the assigned treatment and the received
treatment differ. The assigned-group estimate can understate,
distort, or complicate the effect of actually using the tool.

\item $87/240 = 0.3625$.

\item $120/240 = 0.5$.

\item $66/120 = 0.55$.

\item Not independent because $P(\text{regular}) = 0.3625$, but
$P(\text{regular} \mid \text{open}) = 0.55$.

\item $\binom{5}{2}(0.32)^2(0.68)^3 \approx 0.322$.

\item $E(X) = np = 5(0.32) = 1.6$;
$\text{SD}(X) = \sqrt{np(1-p)} = \sqrt{5(0.32)(0.68)} \approx 1.04$.

\item $(0.68)^3(0.32) \approx 0.101$.

\item Mean $0.50$; SD $\sqrt{0.5 \cdot 0.5 / 120} \approx 0.0456$;
approximately normal because expected successes and failures are
both at least 10.

\item Mean $10$; SD $12/\sqrt{120} \approx 1.095$; approximately
normal by the CLT because $n = 120$ is large.

\item The observed difference is $69/120 - 47/120 = 0.183$. Whether
this is surprising depends on the null sampling distribution.
Later inference shows it is unlikely under equal pass proportions.

\item Larger sample sizes reduce standard error, so the sampling
distribution is narrower. A real difference has a better chance
of producing a large test statistic.

\item $\hat{p} = 69/120 = 0.575$;
$SE = \sqrt{0.575 \cdot 0.425 / 120} \approx 0.0451$;
95\% CI $\approx 0.575 \pm 1.96(0.0451) = (0.486,\ 0.664)$.

\item Difference $= 0.575 - 0.392 = 0.183$;
$SE = \sqrt{\frac{0.575 \cdot 0.425}{120} + \frac{0.392 \cdot 0.608}{120}} \approx 0.0635$;
95\% CI $\approx (0.059,\ 0.308)$.

\item We are 95\% confident the AI-access pass proportion is about
5.9 to 30.8 percentage points higher than the standard-resource
pass proportion.

\item $H_0\!:\ p_{AI} = p_{\text{standard}}$;
$H_a\!:\ p_{AI} > p_{\text{standard}}$.
Pooled proportion $= 116/240 = 0.483$.
$SE = \sqrt{0.483 \cdot 0.517 \cdot (\tfrac{1}{120} + \tfrac{1}{120})} \approx 0.0645$.
$z \approx 2.84$; p-value $\approx 0.0023$; reject $H_0$.

\item If true pass proportions were equal, the probability of getting
a difference at least this large in favor of AI access by random
chance is about 0.0023.

\item Type I: concluding AI access improves pass rate when it does
not. Type II: failing to detect a real improvement from AI access.

\item $18.6 \pm 1.98(12.4/\sqrt{120}) \approx (16.36,\ 20.84)$.

\item Difference $= 8.8$;
$SE = \sqrt{12.4^2/120 + 11.7^2/120} \approx 1.56$;
95\% CI $\approx (5.72,\ 11.88)$.

\item We are 95\% confident AI access increases mean gain by about
5.7 to 11.9 points compared with standard resources.

\item $H_0\!:\ \mu_{AI} = \mu_{\text{standard}}$;
$H_a\!:\ \mu_{AI} > \mu_{\text{standard}}$.
$t \approx 8.8/1.56 = 5.65$; p-value is very small; reject $H_0$.

\item Randomized design makes it more reasonable to attribute the
difference in mean gains to assigned AI access rather than
preexisting group differences.

\item The pretest means are similar: $48.2$ versus $48.9$, so the
data do not support the claim that the AI-access group started
stronger on average.

\item $H_0$: school culture and AI-use category are independent.
$H_a$: school culture and AI-use category are associated.

\item Expected count for restrictive culture and regular AI use:
$120 \cdot 87 / 240 = 43.5$.

\item Expected counts are $43.5$, $36.5$, and $40$ in each row, all
at least 5. The data must also come from an appropriate random
or well-designed study.

\item Degrees of freedom $(2-1)(3-1) = 2$. With $\chi^2 \approx 43.62$,
p-value is less than $0.001$.

\item Reject independence. There is strong evidence of an association
between school culture and AI-use category.

\item The regular-use and no-use cells contribute most. Restrictive
culture has far fewer regular users and far more nonusers than
expected; open culture has far more regular users and far fewer
nonusers than expected.

\item $H_0\!:\ \beta = 0$; $H_a\!:\ \beta > 0$ if testing for
positive association, or $H_a\!:\ \beta \neq 0$ if testing for any
linear association.

\item $t = 0.085 / 0.013 \approx 6.54$.

\item $df = n - 2 = 78$.

\item Reject $H_0$; there is strong evidence of a positive linear
association between AI-tutor minutes and gain score.

\item Using approximately $t^* = 1.99$:
CI $= 0.085 \pm 1.99(0.013) \approx (0.059,\ 0.111)$.

\item Each additional AI-tutor minute per day is associated with an
average gain increase of about $0.059$ to $0.111$ points.

\item Regression here uses actual minutes, which were not randomly
assigned. Motivation and other variables could confound the
relationship.

\item It is too strong because the study directly supports a causal
claim about \emph{assigned access}, not a universal claim that
AI-supported instruction always works. Implementation, student
use, and assessment conditions still matter.

\item It is too strong because misuse or nonuse of a tool does not
prove the tool is ineffective for students who use it as intended.
Noncompliance contaminates the outcome.

\item The intent-to-treat effect compares assigned AI access with
assigned standard resources. The on-treatment association
compares students by what they actually did, such as regular AI
use, but that comparison is more vulnerable to self-selection.

\item A careful conclusion: in this synthetic study, randomized AI
access is associated with higher mean gains and pass rates, and
actual AI use is positively associated with gain. However, actual
use is confounded with student motivation and school culture, so
the cleanest causal conclusion is about access under the study
design.

\item A stronger follow-up study would randomize schools or classes
to clearly defined phone/AI policies, track actual AI use with
logs, require common in-class no-AI assessments, block by school
and teacher, measure prior achievement and motivation, and
compare both intent-to-treat and per-protocol results.
\end{enumerate}

\subsection*{Final Synthesis Guidance}
A strong final answer should say:
\begin{itemize}[leftmargin=*]
  \item The randomized evidence supports a causal claim about
  \textbf{AI access}, not necessarily pure actual AI use.
  \item The actual-use evidence is stronger as an association than as
  a causal claim because of self-selection.
  \item School culture is associated with usage, so local anti-AI or
  pro-AI norms can contaminate interpretation.
  \item AP-style outcomes alone are insufficient if student effort,
  buy-in, or noncompliance differs across groups.
  \item The cleanest conclusion is that AI access, paired with
  independent demonstration requirements, appears to improve
  performance in this synthetic study, but further randomized work
  should separate access, actual use, student motivation, teacher
  implementation, and school culture.
\end{itemize}

\newpage
% ============================================================
\part{Investigative Task: Evaluating an AI-Tutor Pedagogy}
% ============================================================

\begin{scenario}
A high school AP Statistics teacher in Massachusetts implemented an
AI-tutor-supplemented curriculum during the 2024--25 school year. The
class had $n = 22$ students. At year's end, \textbf{all 22 scored a 1}
on the AP exam. The national 1-rate that year was $\mathbf{17\%}$. The
teacher reports that 18 of the 22 students ``boycotted'' the exam
(left it largely blank as protest); the remaining 4 attempted it.

Separately, a randomized controlled trial published in
\emph{Scientific Reports} (Kestin et al., 2025, $N = 194$
undergraduates, randomized crossover) compared an AI tutor to
active-learning classroom instruction on identical content. Reported
result: AI-tutor arm post-test mean $\approx 4.5$, traditional arm
$\approx 3.5$, pooled SD $\approx 1.0$, $p < 10^{-8}$.

The teacher's father, citing the all-1s outcome, argues
\emph{``the data shows your method failed.''}
\end{scenario}

\section{(a) Sampling and external validity \textnormal{\small(Unit 3)}}

The teacher's class is \textbf{not} a random sample from any larger
population. Identify (i) the population the local result can validly
describe and (ii) the population the RCT can validly describe.
Explain why averaging the two as if equally weighted commits a
\textbf{selection bias} error.

\begin{expectedbox}
Local sample $=$ volunteer convenience sample, single classroom,
single teacher in a single Massachusetts district, $n = 22$ ---
describes only itself. RCT $=$ randomized within a defined frame
(Harvard undergrads in physics) --- supports causal inference within
that frame and provides reasonable generalization to similar
learners. The local sample is non-random and confounded with
political/cultural variables; it cannot be pooled as evidence
against the RCT.
\end{expectedbox}

\section{(b) One-proportion z-test on the local outcome \textnormal{\small(Unit 6)}}

Test $H_0: p = 0.17$ vs. $H_a: p > 0.17$ where $p$ is the probability
a student in this class scores a 1.
\[
\hat{p} = 22/22 = 1.00, \qquad
SE_0 = \sqrt{\frac{0.17 \cdot 0.83}{22}} \approx 0.080
\]
\[
z = \frac{1.00 - 0.17}{0.080} \approx 10.4,
\qquad p\text{-value} \approx 0
\]
\textbf{Reject $H_0$.} The 1-rate is wildly higher than national.
Then immediately explain why this rejection does \emph{not} support
the father's causal claim --- go to (c).

\section{(c) Intent-to-treat vs. per-protocol; identifying the estimand \textnormal{\small(Units 3 \& 6)}}

Of 22 students, 18 deliberately did not attempt the exam. Define:
\begin{itemize}[leftmargin=*]
  \item \textbf{Intent-to-treat (ITT) estimand:} mean exam score for
  \emph{all} assigned students. Includes non-compliers.
  \item \textbf{Per-protocol (PP) estimand:} mean exam score for
  students who \emph{actually engaged} with both the instruction and
  the assessment.
\end{itemize}
The all-1s figure is an ITT measure with $\sim$82\% non-compliance.
Argue that the ITT estimate here measures \textbf{student political
compliance with the assessment}, not the causal effect of AI-tutor
instruction on statistical reasoning skill.

\begin{notebox}
\textbf{Analogy:} a drug RCT in which 82\% of the treatment arm
refuses to swallow the pill cannot be used to claim the drug is
ineffective.
\end{notebox}

\textbf{The relevant on-treatment data} are: in-class Blooket
performance, FRQ-style work submitted under observation, verbal
defenses of reasoning. State that the father's argument confuses
ITT non-compliance with causal failure.

\section{(d) Two-sample inference on the RCT \textnormal{\small(Unit 7)}}

Using the published numbers (means $4.5$ vs. $3.5$, pooled
SD $\approx 1.0$, $n_1 = n_2 = 97$):
\[
SE = 1.0 \cdot \sqrt{\tfrac{1}{97} + \tfrac{1}{97}} \approx 0.144
\]
\[
t = \frac{4.5 - 3.5}{0.144} \approx 6.96, \qquad
df \approx 192, \qquad p < 10^{-8}
\]
\textbf{95\% CI for the mean difference:}
$1.0 \pm 1.97(0.144) \approx (0.72,\ 1.28)$.

\textbf{Cohen's $d \approx 1.0$} (very large effect). Conclude: at
standard $\alpha$, we reject $H_0\!: \mu_{AI} = \mu_{\text{trad}}$.
Causal direction is licensed because randomization controls for
confounders by design.

\section{(e) Power and the ``we don't have the data yet'' claim \textnormal{\small(Unit 6)}}

Compute approximate power of the RCT to detect $d = 1.0$ at
$\alpha = 0.05$ with $n_1 = n_2 = 97$:
\[
\text{power} \approx P\!\left(Z > 1.96 - 1.0\sqrt{97/2}\right)
= P(Z > 1.96 - 6.96) \approx 1.000
\]
Power is essentially 1. The claim ``we don't have the data'' is
equivalent to either (i) being unaware of high-powered evidence that
already exists, or (ii) treating $\alpha$ thresholds asymmetrically
(require evidence for AI, accept tradition by default). Both are
inferential errors, not neutral positions.

\section{(f) Confounding in the local sample \textnormal{\small(Unit 3)}}

List confounders in the local AP-score outcome that are \textbf{not}
present in the RCT:
\begin{itemize}[leftmargin=*]
  \item Political identification of students with anti-AI in-group
  \item Teacher reputation effect (admin/peer signaling about the teacher)
  \item Self-selection into the class
  \item Single-teacher / single-school effect (no replication)
  \item Knowledge that the score itself was being used as social currency
\end{itemize}
The RCT randomization breaks the link between treatment and these
variables; the local ``study'' does not. Local outcome is therefore
not a valid test of the treatment effect even before the
non-compliance issue.

\section{(g) Chi-square test for differential response \textnormal{\small(Unit 8)}}

Suppose the teacher records the following $2 \times 2$ table from the
\emph{next} cohort ($n = 30$), classifying students by prior
engagement and post-tutor performance:

\begin{center}
\small
\begin{tabular}{lrrr}
\toprule
 & Improved $\geq$ 1 letter grade & Did not improve & Total \\
\midrule
Previously disengaged & 9 & 3  & 12 \\
Previously engaged    & 4 & 14 & 18 \\
\midrule
\textbf{Total}        & \textbf{13} & \textbf{17} & \textbf{30} \\
\bottomrule
\end{tabular}
\end{center}

Test $H_0$: improvement is independent of prior-engagement status.

Expected counts: e.g., $E_{11} = 12 \cdot 13 / 30 = 5.2$.
Computed $\chi^2 \approx 8.6$, $df = 1$, $p \approx 0.003$.

\textbf{Reject independence.} AI-tutor benefit is \textbf{not
uniform}: marginal gain is concentrated in students the prior system
underserved --- exactly the pattern the tutoring literature predicts
(Bloom's 2-sigma, with largest gains for students far from the
group-paced mean).

\section{(h) Linear regression with interaction \textnormal{\small(Units 2 \& 9)}}

Fit:
\[
\widehat{\text{Score}} = \beta_0 + \beta_1(\text{AI hours})
  + \beta_2(\text{prior engagement})
  + \beta_3(\text{AI hours} \times \text{prior engagement})
\]
Interpret a finding where $\hat\beta_1 > 0$ and $\hat\beta_3 < 0$:
AI hours raise score on average, but the \emph{marginal} benefit
shrinks for already-engaged students. This formalizes the qualitative
claim: \textbf{AI as tutor disproportionately helps students with the
most slack between actual and potential performance}, rather than
uniformly boosting everyone.

\section{(i) Synthesis: stating the defensible conclusion \textnormal{\small(Skills 1.E, 4.B)}}

Write a paragraph, in AP-rubric style, that:
\begin{enumerate}[leftmargin=*]
  \item Identifies the \textbf{causal question} (``Does AI-tutor
  instruction increase statistical reasoning skill relative to
  traditional instruction?'').
  \item Identifies \textbf{which evidence is admissible} for that
  question (RCT --- yes; local ITT outcome under boycott --- no).
  \item States the \textbf{direction and magnitude} of the
  admissible evidence ($d \approx 1.0$, $p < 10^{-8}$).
  \item Identifies the \textbf{failure mode} of the local
  ``evidence'' (selection bias, confounding, near-total
  non-compliance).
  \item Concludes that the empirical claim ---
  \emph{``AI-tutor instruction produces faster skill acquisition
  than traditional instruction, with the largest marginal gains for
  previously underserved students''} --- is \textbf{supported by the
  data available as of 2026}, and that arguments to the contrary
  based on the local outcome rely on inferential errors a Unit 3 /
  Unit 6 student is expected to identify.
\end{enumerate}

\section*{Skills covered}

\begin{center}
\small
\begin{tabular}{cl}
\toprule
Part & AP Stats skill \\
\midrule
(a) & Sampling, external validity, generalizability \\
(b) & One-proportion z-test \\
(c) & Estimand definition, ITT vs PP, interpretation under non-compliance \\
(d) & Two-sample t-inference, CI, effect size \\
(e) & Power, Type II error, asymmetric evidentiary standards \\
(f) & Confounding, randomization as causal license \\
(g) & Chi-square independence, differential treatment effect \\
(h) & Regression, interaction terms, marginal effect interpretation \\
(i) & Scope of inference, synthesis, communication \\
\bottomrule
\end{tabular}
\end{center}

\vspace{1em}
\begin{notebox}
\textbf{Pedagogical point:} a student who can correctly answer
(a)--(i) can reconstruct the argument from first principles, without
needing to share the author's priors. The defense is not ``trust me''
--- it is ``here is the test, and here is what the rubric says.''
\end{notebox}

\end{document}
